\documentclass[11pt]{article}
\usepackage[T1]{fontenc}
\usepackage[utf8]{inputenc}
\usepackage{lmodern}
\usepackage[margin=0.82in]{geometry}
\usepackage{amsmath,amssymb,amsthm,mathtools}
\usepackage{booktabs,longtable,array}
\usepackage{microtype,enumitem}
\usepackage{xcolor}
\usepackage{hyperref,xurl}
\usepackage{fancyhdr}
\hypersetup{colorlinks=true,linkcolor=blue!45!black,citecolor=blue!45!black,
urlcolor=blue!45!black,pdftitle={Signed horizons and residue-encoded Pell witnesses},
pdfauthor={AI-assisted research prepared for Vladimir Reshetnikov}}
\setlength{\headheight}{14pt}
\pagestyle{fancy}
\fancyhf{}
\fancyhead[L]{\small Explicit Diophantine forms of consistency}
\fancyhead[R]{\small 25 October 2026}
\fancyfoot[C]{\thepage}
\setlength{\emergencystretch}{2em}
\setlist{nosep}
\allowdisplaybreaks[3]
\newtheorem{theorem}{Theorem}[section]
\newtheorem{lemma}[theorem]{Lemma}
\newtheorem{proposition}[theorem]{Proposition}
\newtheorem{corollary}[theorem]{Corollary}
\newcommand{\Z}{\mathbb Z}
\newcommand{\N}{\mathbb N}
\newcommand{\Con}{\operatorname{Con}}
\newcommand{\Refut}{\operatorname{Ref}}
\newcommand{\PellX}{\chi}
\newcommand{\PellY}{\psi}
\title{\textbf{Signed horizons and residue-encoded Pell witnesses}\\[5pt]
\large Explicit Diophantine equations for the consistency of ZFC}
\author{Research prepared for Vladimir Reshetnikov and ProveIt\\
\small AI-assisted mathematical development with separate proof and artifact reviews}
\date{25 October 2026}
\begin{document}
\maketitle
\begin{abstract}
We give new explicitly instantiated Diophantine equations whose
unsolvability, in their stated domains, is provably equivalent in ZFC
to its syntactic consistency. A signed affine proof horizon removes an
integer variable from each of two directly compiled quartics. The balanced
result has 214 unrestricted integer variables, 153 squared quadratic
residuals, 5,875 collected monomials and coefficient height 655,360.
A second quartic has 212 integer variables and larger coefficients.
An independent Pell-ordinate construction replaces a quantified quotient
by a residue already carried by an auxiliary ordinate. It yields fourteen
positive variables at 66 operations and degree 174, or 68 operations and
degree 146. A careful domain audit turns the latter into a fully explicit
equation in 36 unrestricted integers, degree 288 and 145 operations.
Composing signed horizons with a reused CRT anchor and oriented power
pooling also gives eight direct integer profiles, including 117 variables
at degree 16, 108 at degree 32, and 86 at degree 2,097,780.
All final equations have fixed coefficients and contain only polynomial
operations. The new proofs are ordinary mathematics formalizable in ZFC,
composed with explicitly identified inherited calibration theorems.
Exact finite verification is supplied, but no new end-to-end
Lean/Rocq or encoded object-language ZFC proof is claimed.
\end{abstract}

\newpage
\tableofcontents
\newpage
\section{Results, interpretation, and comparison scope}
\label{sec:results}

The task is to make the consistency equation simpler while keeping both
directions of its calibration. There is no single order on simplicity:
variable count, total degree, evaluation length, coefficient height,
coefficient storage and proof complexity can move in opposite directions.
This submission therefore gives exact completed constructions and makes
the relevant tradeoffs explicit. It does not prove global optimality or
claim worldwide priority for its component techniques.

The repository comparison is pinned to
\[
\texttt{02a78a08603cbae248c19c23a75907c6c6bd6443}.
\]
The immediate direct source is package 173~\cite{proveit173}; the compact
source is package 174~\cite{proveit174}, with a residue-carrier argument
from package 170~\cite{proveit170}. The direct higher-degree family also
uses the separately delivered \emph{Oriented Powers} packet~\cite{oriented}.
Its archive identity is retained in the source ledgers.

\subsection{The precise logical claim}

Fix the effective first-order presentation denoted \texttt{ZFCax\_s}
and the finite proof relation \texttt{BProv} in the repository.
Let $\Refut_{\mathrm{ZFC}}$ mean that there is a finite derivation of
\texttt{Form.fBot} from that presentation, and put
$\Con(\mathrm{ZFC})=\neg\Refut_{\mathrm{ZFC}}$.
The new constructions prove the following ordinary theorem.

\begin{theorem}\label{thm:main}
For every final polynomial $P$ delivered with this article and its
specified domain $\mathcal D$, ZFC proves
\[
 \bigl(\exists\mathbf x\in\mathcal D\;P(\mathbf x)=0\bigr)
 \longleftrightarrow \Refut_{\mathrm{ZFC}}.
\]
Consequently ZFC proves
\[
 \boxed{\bigl(\neg\exists\mathbf x\in\mathcal D\;P(\mathbf x)=0\bigr)
 \longleftrightarrow \Con(\mathrm{ZFC}).}
\]
The coefficients are fixed independently of the truth of consistency.
\end{theorem}

The theorem proves an equivalence. It does not establish the absence of
zeros, nor consistency itself. Its quantifiers are the internally coded
integers and finite derivations of ZFC. The argument is not restricted to
externally standard proof lengths.

\subsection{The useful completed endpoints}

\begin{center}
\small
\begin{tabular}{lrrrr}
\toprule
Equation family & Domain & Variables & Exact degree & Binary operations\\
\midrule
Balanced direct quartic $Q_{214}$ & $\Z$ &214&4&---\\
Dimension-focused quartic $Q_{212}$ & $\Z$ &212&4&---\\
Compact $P_E$ & $\Z_{>0}$ &14&174&66\\
Compact $P_q$ & $\Z_{>0}$ &14&158&67\\
Compact $P_q^\circ$ & $\Z_{>0}$ &14&146&68\\
Explicit compact integer lift $P_{36}$ & $\Z$ &36&288&145\\
Direct pooled profile & $\Z$ &132&12&2484\\
Direct pooled profile & $\Z$ &117&16&2439\\
Direct pooled profile & $\Z$ &108&32&2434\\
Smallest delivered direct profile & $\Z$ &86&2097780&2385\\
\bottomrule
\end{tabular}
\end{center}

Operation counts are binary $+,-,\times$ gates in the supplied
straight-line programs, with fixed coefficients as input numerals
and shared intermediate values counted once. They measure arithmetic
circuit length; big-integer bit complexity is a separate cost.

The compact integer equation $P_{36}$ has a short factorized definition:
Section~\ref{sec:ordinate} gives all its polynomial factors and
Section~\ref{sec:ordinate-constants} defines its four fixed constants.
The balanced quartic has modest \emph{expanded} coefficients and no
native coefficient ports; Appendix~\ref{app:quartic} prints every residual.
The direct circuits are useful alternatives when avoiding the large
compact constants matters more than arithmetic evaluation length.

The fourteen-positive-variable equations do not assert a
fourteen-\emph{integer}-variable theorem. The stated conversion to 36
integer variables is part of the proof. Likewise, the 86-variable direct
equation is not a variable-count record across other encoding families.
Much smaller-variable equations at other degree and coefficient costs
already appear in the repository. Published small-variable universality
results also have a different optimization objective~\cite{sun2026}.

\subsection{How this iteration changes the earlier results}

The balanced parent from package 173 has 215 integer variables, 154
quadratic residuals, 5,660 monomials and height 91,679. The new balanced
quartic removes one variable and one residual, at the cost of increasing
support to 5,875 and height to 655,360. The 213-variable companion
similarly becomes 212 variables, with increased coefficient height.
These are dimension reductions, not improvements in every metric.

The compact parent has fifteen positive variables, degree 157 and 66
operations. The new $P_E$ keeps 66 operations with fourteen variables,
while its degree is 174. The lower-degree $P_q^\circ$ improves the
documented fourteen-positive 71-operation, degree-157 comparison to
68 operations and degree 146. Complete coefficient tuples are kept
separate throughout this comparison.

The pooled direct family has stronger matched improvements:
137 variables at degree 12 become 132 with the same bound of 256 on
literal circuit constants; 121 at degree 16 become 117; 111 at degree 32
become 108. Its final 86-variable profile has degree 2,097,780, compared
with 87 variables and degree 8,388,612 in the earlier packet. The
expanded coefficient bounds remain a separate column of the exact
machine-readable comparison.

\section{A signed horizon removes a quartic variable}
\label{sec:quartic}

The old direct source imposed a positive quadratic horizon such as
$t=16+4z^2$ or $t=128+4z^2$. This guaranteed positivity at the cost of
a coordinate that could not be removed affinely. The power machinery
already contains enough information to recover the sign. We use
that information and replace the square progression by a signed
affine progression.

\subsection{A uniform signed-horizon lemma}

Put $\ell=2^{20}$. Fix an integer $m\ge8$ divisible by four and impose
\[
 t=m(2z+1),\qquad z\in\Z.
\]
Then $t\ne0$, $|t|\ge m$ and $4\mid t$. Retain the source's signed
Pell frame, its reused even CRT index $C$, all orientation rows and
all actual power requests. Two growth choices will be used:
\begin{align}
 K&=\ell t, \label{eq:linear-growth}\\
 U&=a t^2,\qquad K=bU^2
 \quad(a,b\in\Z_{>0},\;ba^2m^3\ge\ell).
 \label{eq:quartic-growth}
\end{align}
The retained Pell request expresses $F=t^K$. This exponentiation
is part of the \emph{proof interpretation}, not the final polynomial.

\begin{lemma}\label{lem:signed-horizon}
At every integer zero of either modified source, $t\ge m$ and
$K\ge\ell t$. Every finite C222 refutation can be represented using
the new positive branch of the horizon progression.
\end{lemma}

\begin{proof}
The source's preliminary signed anchor argument uses squares and
absolute values. It recovers the canonical signed Pell pairs before
decoding arithmetic signs. Its actual frame includes the squares of
all requested bases, outputs and indices, in particular the changed
index $K$. Thus it still gives
\[
 W\ge258,\qquad |A|>W^W.
\]
For the affine index this lower bound on $W$ already follows from
$K^2\ge\ell^2m^2$; for the quartic index it follows from
$K=ba^2t^4$. The signed power lemma therefore applies at a nonzero
base. At a base of absolute value at least two, an integral exact
signed-index power must have nonnegative exponent.

For \eqref{eq:linear-growth}, $|t|\ge8$ makes the first power
force $K\ge0$. Since $K=\ell t$ and $t\ne0$, this proves $t\ge m$
immediately.

For \eqref{eq:quartic-growth}, the index $K=ba^2t^4$ is positive
and even before the sign of $t$ is known. The first power gives
$F=t^K>0$, with $F$ even and $F>|t|$. The retained CRT row is
\[
 Bj=C+F+1,\qquad C=2c.
\]
Its right side is odd, so $B$ is odd and nonzero. The self-radix
multiplier has the form $H=2+4S$ for an integer $S$, hence
$|BH|\ge2$. The exact power at that base gives
$x\ge0$ and $q=(BH)^x$.

It remains to rule out $x=0$. The beta endpoint rows, which were
never deleted, give
\[
 y=F(2v+1),\qquad g=1+(t+1)y,\qquad x=1+gu_t.
\]
Because $|y|\ge F$ and $|t+1|\ge m-1\ge7$,
\[
 |g|\ge(m-1)F-1>1.
\]
Thus $x=0$ would imply $gu_t=-1$, impossible in integers.
We have $x\ge1$ and consequently $|q|\ge2$. The literal retained
request is $q^t$. Its signed exactness forces $t\ge0$, and the
progression then gives $t\ge m$. Notice that the proof does not
replace the digit row by a statement about $B$ alone; its unrestricted
high quotient would invalidate that shortcut.

For either growth choice, positivity now yields $K\ge\ell t$.
Completeness uses the cofinal positive progression $m(2z+1)$:
choose $t$ above the finite proof length and the local free parameters,
then use the source construction with this $t$. The numerical
estimates and the order of the subsequent choices are justified below.
\end{proof}

\subsection{The estimates needed at minimum eight}

The literal square-shell constructors obey a sharper bound than
the earlier conservative $t\ge16$ presentation. For integers
$0\le r,s<T$,
\[
 E(r,s)\le8T^2-2T+1<8T^2,
\]
and the other three constructors are no larger. Thus, if both
arguments are strictly below $t^e$ and $t\ge8$, every constructor
output is strictly below $t^{2e+1}$. With
$e_0=1$ and $e_{d+1}=2e_d+1$, one obtains
$e_d=2^{d+1}-1$.

The authenticated source constructor depths are
\[
 4,\;3,\;3,\;3,\;3,\;5,\;4,\;3,\;19.
\]
Consequently every local formula label is below
$t^{2^{20}-1}$. Set $L=t^\ell$. Then $L\ge t$ and
$F=t^K\ge L^t$. No bound on the deleted sum of derived labels is
needed here.

\begin{lemma}[Boundary-valid label separation]\label{lem:label}
Let $t,L,F$ be positive integers with $L\ge t$ and $F\ge L^t$.
If two multisets of labels in $[1,L]$, with at most $t$ and
$2t+1$ entries respectively, satisfy
\[
 \prod_{j\in U}(1+jF)=\prod_{j\in V}(1+jF),
\]
then their label multiplicities agree.
\end{lemma}

\begin{proof}
If the factor counts differ, let $r$ be the smaller count.
Then $r\le t$, and comparison of the two products gives
\[
 (F+1)^{r+1}\le(1+LF)^r\le L^r(F+1)^r,
\]
contradicting $F+1>L^t$. Their common count $s$ is therefore
at most $t$. If $s=0$, both multisets are empty; assume henceforth
$s\ge1$. Each coefficient of either product polynomial
in an indeterminate lies between $1$ and
\[
 \binom{s}{k}L^k\le s^{s-k}L^k\le L^s\le F.
\]
Thus every coefficient of their difference has absolute value
strictly below $F$, even when one coefficient equals $F$.
The least nonzero coefficient of a difference vanishing at $F$
would have to be divisible by $F$, a contradiction.
The product polynomials agree. Their roots $-1/j$ recover
the multiplicities.
\end{proof}

Truth has label zero and hence weight $1+0F=1$; its factors are
omitted before applying the positive-label lemma.
The source's multiplicative Horn updates therefore recover the
same finite rule-flow balance. If falsity were outside the closure
of truth under the selected rules, every outside conclusion would
require an outside premise. Summing that balance over outside labels
would be nonpositive, whereas the target falsity contributes one.
This yields a finite C222 refutation even if the selected rules have
cycles.

Two other conservative numerical estimates have explicit replacements.
Since $K\ge\ell t$ and $t\ge8$,
\[
 F\ge8^{8\ell}>2^{64}.
\]
After the unchanged beta sign recovery, $y\ge F$ and
$x\ge2+(t+1)y$. Therefore
$2+t+x+y+F<2x$, $x>2^{32}$ and the retained local polynomial budget
gives
\[
 \mathrm{Local}<(2x)^{2^{23}}<2^x.
\]
Also
\[
 F\ge8^{2t}>16^t>
   \operatorname{lcm}(1,\ldots,t)^2.
\]
This supplies the radius for the inherited factorial-free binomial
CRT selector. No factorial-rounding lemma is being applied outside
its range.

Finally
$v_2(F)=K v_2(t)\ge\ell t$ exceeds both
$\lfloor\log_2t\rfloor$ and $v_2(t!)$. Hence
$\binom{F+t}{t}$ is odd. Multiplying $F$ by a suitable positive
odd multiple supplies a beta spacing divisible by $t!$ and larger
than all finite pending states. After choosing the beta code and
the every-prime CRT residues, choose the encoded index $C$ even
and larger than all other requested indices. Only then choose
$H,q$ and the Pell frame. The choices are cofinal and noncircular.
These facts complete the soundness and completeness assertions
in Lemma~\ref{lem:signed-horizon}.

\subsection{Two literal quartics}

The dimension-focused choice is
\[
 t=8+16z,\qquad K=1048576t.
\]
Affine elimination of $t$ gives $Q_{212}$ with 212 integer
coordinates, 151 quadratic residuals and 5,706 collected monomials.
Its exact expanded coefficient height is
\[
 316912650058804747029598371840.
\]

The balanced choice is
\[
 t=8+16z,\qquad U=2t^2,\qquad K=512U^2=2048t^4.
\]
At $t=8$ the inequality $2048t^3\ge2^{20}$ is already an equality.
This gives $Q_{214}$ with 214 integer coordinates, 153 residuals,
5,875 collected monomials and exact coefficient height 655,360.
Its maximum residual coefficient is 512.

For each profile, replace the source's horizon row by
$t-m-2mz=0$, substitute $t=m+2mz$ into every other row, and
delete this row and the coordinate $t$. This is a polynomial
bijection on zeros of the \emph{modified} source, with inverse
restoring the displayed $t$. Affine substitution keeps every
residual quadratic.

For the balanced profile, change only the two protected growth
rows to $U-2t^2$ and $K-512U^2$ before eliminating $t$.
Subtract twice the square-definition row from the frame row,
and add the square-definition row to the orientation row.
These integral row operations have explicit inverses and remove
duplicated $t^2$ terms.

The finite checker also verifies that these edits commute with
the inherited six coordinate eliminations and optimization chain.
The radical directions vanish on the protected coordinates;
the affine projection uses only the stated local constructor
rows; later charts leave the horizon and growth coordinates
unchanged; and row operations do not modify their defining rows.
Thus the emitted equations implement the stated arithmetic
source change, rather than an unsupported edit to an optimized
polynomial.

\begin{corollary}
For $r=212,214$, ZFC proves
\[
 \neg\exists\mathbf x\in\Z^r\;Q_r(\mathbf x)=0
 \quad\longleftrightarrow\quad\Con(\mathrm{ZFC}).
\]
\end{corollary}
The complete $Q_{214}$ residual list appears in
Appendix~\ref{app:quartic}; both collected polynomials are in
\texttt{quartic/}.

\section{An auxiliary ordinate removes a quantified quotient}\label{sec:ordinate}

The next construction improves the inexpensive positive-variable branch.
It combines the normalized auxiliary ordinate of the fifteen-variable
P66 equation~\cite{proveit174} with the residue-carrier argument developed
for combined Pell roots in package~170~\cite{proveit170}. The resulting
equations have fourteen positive unknowns and the following exact profiles.
The whole Start-only4127 coefficient tuple is retained unchanged; its
complete finite definition is in Section~\ref{sec:ordinate-constants}.

\begin{center}
\begin{tabular}{lrrr}
\toprule
Equation & Positive unknowns & Binary operations & Degree\\
\midrule
Immediate P66 parent &15&66&157\\
Previously documented comparison &14&70&189\\
Previously documented comparison &14&71&157\\
New $P_E$ &14&66&174\\
New $P_q$ &14&67&158\\
New $P_q^{\circ}$ &14&68&146\\
\bottomrule
\end{tabular}
\end{center}

The first new row reduces the dimension of the immediate parent at the same
arithmetic cost, with a degree increase. The first and last new rows also
improve both cost and degree against the corresponding fourteen-variable
comparisons. The previous comparison tuples and the new tuple are distinct
complete coefficient specializations. These are documented improvements
within this branch, not a global minimum theorem.

\subsection{Complete polynomial definitions}

The independent positive integers are
\[
 (J,F,\alpha,z,f,n,i,T,s,w,\tau,\eta,\zeta,\gamma).
\]
With the four fixed integers $\beta,K_0,M_C,M_F$, define
\begin{align*}
q&=\beta J+1,&C&=q-F-\alpha,&X&=wq,&Y&=sq^3,\\
E&=XY,&k&=\eta+\zeta,&\ell&=EkY,&c&=kY+\eta,\\
a&=E+Y,&H&=4a+3,&\Delta&=a^2+H=(a+2)^2-1,
\end{align*}
and
\begin{align*}
D&=X+ac+\gamma H,\\
R&=[(q-1)(q-F)+\alpha](q^2-1)+(M_C+qM_F)J,\\
V&=f^2(cT-R)-c,\\
N_0&=\tau^2-\ell(\ell+k),\qquad N_1=D^2-\Delta c^2,\\
N_t&=(K_0+w)C+(q-F)-z(q-1).
\end{align*}
Choose either $(M,d)=(E,E)$ or $(M,d)=(q^4,q^2)$ and put
\begin{align*}
u&=k-Mn-R,&I&=dic^2,&S&=\Delta I,\\
N_a&=S^2(V^2-u^2)+u^2,\\
N_s^{\mathrm{sc}}&=\Delta f^2-S^2,&N_s&=f^2-\Delta I^2.
\end{align*}
The displayed letters are polynomial abbreviations, not additional
unknowns. For $(M,d)=(E,E)$ the equation is
\begin{equation}\label{eq:ordinate66}
 \boxed{P_E=N_0N_1N_aN_tN_s^{\mathrm{sc}}-\Delta=0.}
\end{equation}
For $(M,d)=(q^4,q^2)$ we supply both
\begin{equation}\label{eq:ordinate67}
 P_q=N_0N_1N_aN_tN_s^{\mathrm{sc}}-\Delta=0.
\end{equation}
\begin{equation}\label{eq:ordinate68}
 \boxed{P_q^{\circ}=N_0N_1N_aN_tN_s-1=0.}
\end{equation}
The exact identity $P_q=\Delta P_q^{\circ}$ and $\Delta>0$ show that
these last two equations have the same positive zero set. The unscaled
strong norm costs one additional multiplication but reduces degree by twelve.

\paragraph{The equivalence theorem.}
For each of these fixed polynomials, ZFC proves that a positive zero exists
if and only if a finite syntactic ZFC refutation exists. Consequently ZFC
proves that the stated positive zero set is empty if and only if
$\operatorname{Con}(\mathrm{ZFC})$.
The proof below restores a literal parent P66 zero and constructs new
witnesses from its canonical completeness data. The native and foundational
calibrations are the identified inherited theorems, not conclusions of a
finite numerical search.

\subsection{Unit signs are resolved before recovering an index}

At a positive zero, canceling $\Delta$ in the proof, if necessary, gives
\begin{equation}\label{eq:ordinateunits}
 N_0N_1N_aN_tN_s=1.
\end{equation}
Every factor is an integer unit. Let $v=XY^2=EY$. Then
$N_0=\tau^2-v(v+1)k^2$. For $v\ge2$, the negative Pell equation
$x^2-v(v+1)y^2=-1$ is impossible: a solution with least positive $y$
satisfies $vy<x<(v+1)y$, and multiplication by
$(2v+1)-2\sqrt{v(v+1)}$ produces another solution with ordinate
$0<|(2v+1)y-2x|<y$. Thus $N_0=1$.

The number $\Delta=(a+2)^2-1$ is $0$ or $3$ modulo $4$, so neither
$N_1$ nor $N_s$ can be $-1$. Modulo $4$, $N_a$ is congruent to $u^2$
or $V^2$, according as $S$ is even or odd, so it cannot be $-1$ either.
Equation~\eqref{eq:ordinateunits} therefore yields
\begin{equation}\label{eq:ordinateallunits}
 N_0=N_1=N_a=N_t=N_s=1.
\end{equation}
In particular, deleting the old quotient factor removes its previously
unresolved partner sign in $N_t$. We immediately have
\[
 (K_0+w)C=1+z(q-1)-(q-F)\ge F>0,
\]
so $C>0$. The unchanged coherent coefficient tuple satisfies
$0<M_C<\beta<M_F<2\beta$. Applying these inequalities to the literal
definition of $R$, with
$G_0=q\alpha+(q-1)C$, gives
\begin{equation}\label{eq:ordinatemask}
 3q<R<q^4-q^3+q^2-1<q^4-2\le E-2.
\end{equation}
In both variants $R<M$. None of these statements has used the omitted
quotient or assumed a main Pell index.

\subsection{The main index is authenticated without that quotient}

For $b\ge2$, write $\chi_b(j),\psi_b(j)$ for the Pell sequences with
initial values $(1,b)$ and $(0,1)$ and recurrence
$U_{j+2}=2bU_{j+1}-U_j$. Put $A=a+2$ and $P_0=1+2v$.
The norm equations imply
\begin{align}
 \tau&=\chi_{P_0}(r),&k&=2\psi_{P_0}(r),&r&\ge1,
     \label{eq:ordinatefirst}\\
 D&=\chi_A(p),&c&=\psi_A(p),&p&\ge1.
     \label{eq:ordinatemain}
\end{align}
For the first norm the fundamental positive solution is $(2v+1,2)$;
ordinate one would put a square strictly between $v^2$ and $(v+1)^2$.
The case $r=1$ would give $k=2$, $\eta=\zeta=1$, and
$1<c=2Y+1<2A$, which is impossible in~\eqref{eq:ordinatemain}.
Hence
\begin{equation}\label{eq:ordinategrowth}
 r\ge2,\quad k\ge8XY^2+4,\quad c>kY>8XY^3>2R.
\end{equation}
Since $P_0>A$ and $c>k$, monotonicity gives $r<p$; now $p\ge3$ and
$A\ge2$ imply $c=\psi_A(p)>2p$.

The strong equation gives $f=\chi_A(m)$ and $I=\psi_A(m)$.
The literal divisibility $c^2\mid I$ and strong divisibility first imply
$p\mid m$. Writing $m=pj$, composition gives
$I=c\psi_D(j)$. Since $D^2\equiv1\pmod c$, the recurrence gives
$\psi_D(j)\equiv jD^{j-1}\pmod c$. Thus $c\mid j$ and
\begin{equation}\label{eq:ordinatestrongindex}
 pc\mid m,\qquad c\mid m,\qquad m>2p.
\end{equation}

Let $Q=S^2=\Delta(f^2-1)>1$. The auxiliary equation is
\[
 QV^2-(Q-1)u^2=1.
\]
Neither coordinate is zero; also $V>0$ by
\eqref{eq:ordinatemask}--\eqref{eq:ordinategrowth}, $T\ge1$, and $f^2>2$.
Therefore $SV=\chi_S(\nu)$ and $|u|=\psi_S(\nu)$ for a positive odd
index $\nu$. An even index would make the abscissa $\pm1$ modulo $S$.

For $\nu=2j+1$, the integer polynomials
$L_0(Z)=1$, $L_1(Z)=4Z-3$,
$L_{j+1}(Z)=(4Z-2)L_j(Z)-L_{j-1}(Z)$ satisfy
\begin{align*}
 \chi_b(2j+1)&=bL_j(b^2),&L_j(0)&=(-1)^j(2j+1),\\
 L_j(1-A^2)&=(-1)^j\psi_A(2j+1).
\end{align*}
Using $Q\equiv1-A^2\pmod f$ and $V\equiv-c\pmod f$ yields
$\psi_A(\nu)\equiv\pm c\pmod f$. Squaring gives
\[
 \chi_A(2\nu)\equiv\chi_A(2p)\pmod{\chi_A(m)}.
\]
The signed step-down argument from~\cite{proveit170} applies with
$0<2p<m$: addition identities reduce any abscissa to a signed
$\chi_A(t)$ with $0\le t\le m$; the endpoint has residue zero,
$2\chi_A(m-1)<\chi_A(m)$ excludes a negative representative, and
monotonicity identifies the positive one. It follows that
$2\nu\equiv\pm2p\pmod{2m}$ and hence $\nu\equiv\pm p\pmod c$.
Separately, $Q\equiv0\pmod c$ and $V\equiv-R\pmod c$ give
$\nu\equiv\pm R\pmod c$. As $0<p,R<c/2$, we conclude
\begin{equation}\label{eq:ordinatepR}
 p=R.
\end{equation}
The omitted quotient has not occurred in any premise of this authentication.

\subsection{The ordinate restores the missing quotient}

Both choices satisfy $M\mid v$ and $M\mid S^2$. Reducing the first
recurrence modulo $M$ gives $k\equiv2r\pmod M$. The odd-ordinate identity
\begin{equation}\label{eq:ordinateodd}
 \psi_S(2j+1)\equiv(-1)^j\pmod{4S^2}
\end{equation}
follows by induction. Thus $u\equiv\pm1\pmod M$. Since
$u=k-Mn-R$, there is $\varepsilon\in\{1,-1\}$ such that
$2r-R\equiv\varepsilon\pmod M$. Now $r<p=R<M$ implies
$|2r-R-\varepsilon|<R<M$, so
\[
 2r=R+\varepsilon.
\]
The negative choice would give $R=2r+1$. But $B_0=2A^2-1>P_0$ and
composition then imply
\[
 c=\psi_A(2r+1)>2A\psi_{B_0}(r)
 \ge2A\psi_{P_0}(r)=Ak>k(Y+1),
\]
contrary to $c=kY+\eta<k(Y+1)$. Therefore $2r=R+1$.
The full first recurrence gives $k\equiv2r\pmod{4v}$, so
\begin{equation}\label{eq:ordinatehrestore}
 h_{\rm old}=\frac{k-R-1}{E}\in\mathbb Z_{>0}.
\end{equation}
Positivity follows from~\eqref{eq:ordinategrowth} and~\eqref{eq:ordinatemask}.
Together with $i_{\rm old}=di$ and $y_{\rm old}=|u|>0$, this restores
an exact zero of the original P66 equation: its omitted quotient factor
is $+1$, and all its other normalized factors are $+1$.
The retained whole-tuple calibration supplies a finite syntactic refutation.

For completeness, start with a refutation and its canonical outer P66
witnesses. The inherited native theorem supplies $R\equiv3\pmod4$,
$c=\psi_A(R)$, $k=2\psi_{P_0}((R+1)/2)$, and the required outer unit
factors. Choose a positive period of the determinant-one Pell matrix
modulo $c^2d$. It supplies $m>0$ such that $c^2d\mid\psi_A(m)$.
Set $f=\chi_A(m)$ and $i=\psi_A(m)/(c^2d)>0$ before choosing fresh
auxiliary witnesses. Define
\[
 V_* = \frac{\chi_S(R)}S,\qquad y_* =\psi_S(R).
\]
These are positive integers. The same polynomial identities, now modulo
$f^2$ and $c$, give $V_*\equiv-c\pmod{f^2}$ and
$V_*\equiv-R\pmod c$. Since $\gcd(c,f)=1$ and $f^2\equiv1\pmod c$,
\[
 T=\frac{V_*+c+Rf^2}{cf^2}\in\mathbb Z_{>0}.
\]
Finally~\eqref{eq:ordinateodd} gives $y_*\equiv-1\pmod M$, whereas
$k-R\equiv1\pmod M$. Thus
\[
 n=\frac{y_*+k-R}{M}\in\mathbb Z_{>0},\qquad u=-y_*.
\]
Every new coordinate is positive and every unit factor is $+1$. This proves
completeness with newly chosen private witnesses, rather than assuming old
private witnesses retain stronger divisibilities.

\subsection{Exact costs and nonzero leading forms}

In $P_E$ the parent's three gates for $k-hE-R$ are reused as $k-nE-R$.
Removing this factor from the final product saves one multiplication;
scaling $S$ by $E$ costs one. Thus the count remains 66. The $q^4$ modulus
adds one gate, giving 67; computing $SI$ for the unscaled strong norm adds
one further gate, giving 68. All inputs and gates are live.

Write $q_0=\beta J$, $k_0=\eta+\zeta$ and
$B_0=w(q_0-F-\alpha)-zq_0$. The complete leading homogeneous forms are
\begin{align*}
 [P_E]_{174}&=8q_0^{103}\gamma i^4k_0^{13}w^{16}s^{29}T^2f^4B_0,\\
 [P_q]_{158}&=8q_0^{95}\gamma i^4k_0^{13}w^{12}s^{25}T^2f^4B_0,\\
 [P_q^{\circ}]_{146}&=8q_0^{87}\gamma i^4k_0^{13}w^{10}s^{23}T^2f^4B_0.
\end{align*}
The only leading cancellation is accounted for exactly by
\[
 N_1=2acZ+Z^2-Hc^2,\qquad Z=X+\gamma H,
\]
which has degree 18. The five factor degrees are respectively
$(22,18,74,2,58)$, $(22,18,66,2,50)$ and $(22,18,66,2,38)$.
Choosing the $-zq_0$ term and the $\eta^{13}$ term displays coefficients
$-8\beta^{104}$, $-8\beta^{96}$ and $-8\beta^{88}$. They are nonzero at
the actual fixed $\beta$, certifying the claimed exact degrees.

\subsection{A complete equation in 36 unrestricted integers}\label{sec:ordinate36}

The sign audit permits three parent coordinates to range over all integers.
The polynomial uses $f$ and $\tau$ only squared, and their recovered norms
force them to be nonzero, so absolute values restore positivity.
The sign of $n$ was never used in soundness: only the signed carrier $u$
and its residue modulo $M$ matter. Completeness already constructs these
three coordinates positively. Thus only
\[
 J,F,\alpha,z,i,T,s,w,\eta,\zeta,\gamma
\]
require positivity encodings.

For each coordinate $x$ in this eleven-element list, substitute
\begin{equation}\label{eq:ordinatepositiveparameter}
 x=1+x_a^2+x_b^2+x_c(x_c+1),\qquad x_a,x_b,x_c\in\mathbb Z,
\end{equation}
in $P_q^{\circ}$ and leave $f,n,\tau$ unrestricted. This is the explicit
polynomial $P_{36}$ in the supplied literal graph. The parameterization is
positive for every integer triple and represents every positive integer.
For a short derivation, Legendre's three-square theorem applies to $4N+1$
for every $N\ge0$; reducing modulo 4 shows that exactly one square
is odd, so after permutation
\[
 4N+1=(2a)^2+(2b)^2+(2c+1)^2,
 \quad N=a^2+b^2+c(c+1).
\]
This representation is also explicitly formalized in the Archive of Formal
Proofs~\cite{danilkin2023}. We obtain the ordinary ZFC theorem
\begin{equation}\label{eq:ordinateintegercon}
 \boxed{\mathrm{ZFC}\vdash
  \bigl(\neg\exists\mathbf x\in\mathbb Z^{36}\;P_{36}(\mathbf x)=0\bigr)
  \ \longleftrightarrow\ \operatorname{Con}(\mathrm{ZFC}).}
\end{equation}
No inequalities or extra quantified registers remain in $P_{36}$.

There are $3\cdot11+3=36$ integer unknowns. Each substitution costs seven
binary operations, giving $68+7\cdot11=145$ operations. Giving the eleven
substituted coordinates weight two and $f,n,\tau$ weight one yields factor
degrees $(44,36,128,4,76)$, totaling 288. The squared carrier contributes
only degree 118 in the auxiliary norm, so it cannot overtake the degree-128
$V^2$ term. The complete leading form is the last displayed positive
leader with each substituted $x$ replaced by
$\sigma_x=x_a^2+x_b^2+x_c^2$ and with $f$ unchanged. Setting all $x_b,x_c$
to zero exhibits the nonzero coefficient $-8\beta^{88}$ at
\[
 J_a^{176}\gamma_a^2 i_a^8\eta_a^{26}w_a^{20}s_a^{46}T_a^4f^4z_a^2.
\]
Its exponents sum to 288, so this degree is exact.

The 36-integer equation uses a large fixed coefficient architecture. Its
size should therefore be compared separately with small-coefficient direct
quartics and with still smaller but higher-degree constructions. The new
finite checks authenticate the literal formulas, full fixed bytes, liveness,
and leading forms; the unbounded equivalence is the mathematical proof above
composed with the retained native/refutation calibration.

\section{Composing signed horizons with oriented power pools}

The reused CRT anchor and the oriented power theorem have complementary
strengths. The former removes a derived-label range request; the latter
combines the remaining power certificates. Their composition requires a new
sign argument, because the factorial normalization used in the earlier
pooling proof is absent after the anchor reduction.

\subsection{The common arithmetic interface}

The five bases are
\[
 n_0=t,\qquad n_1=B,\qquad n_2=BH,\qquad n_3=q,\qquad n_4=q+1,
\]
with respectively $1,1,1,3,10$ power requests. The retained canonical signed
Pell data at center $A$ satisfy
\[
 |A|>W^W,\qquad n_j^2,m_i^2,k_i^2<W,\qquad W\ge258.
\]
This interface is established before the signs of the arithmetic coordinates
are decoded. Every requested index has a canonical signed Pell pair
$(X_i,Y_i)$, with $D_i=X_i-AY_i$.

Define, for each block,
\begin{align*}
 T_j&=2An_j-n_j^2-1,\\
 E_j&=\sum_{i\in j}\bigl((m_i-D_i)^2-Y_i^2\bigr)
       -2n_j\sum_{i\in j}Y_i(m_i-X_i).
\end{align*}
The five-base residual is the ordinary polynomial
\[
 R=\sum_{j=0}^4 n_jE_j\prod_{\ell\ne j}T_\ell
       -z\prod_{j=0}^4T_j.
\]
The inherited oriented-pooling theorem makes $R=0$ equivalent to exactness
of every request at a nonzero actual base. Zero-base blocks are temporarily
unconstrained, and are excluded by the argument below. For nonzero bases,
$\operatorname{sgn}(n_j/T_j)=\operatorname{sgn}(A)$; this common sign is what
prevents cancellation of the squared power errors.

Two smaller-degree alternatives use three and two quotient coordinates.
The paired version retains the direct $t$ singleton, pools $B,BH$ with base
weights, and pools $q,q+1$ without weights. The staged version pools
$t,B,BH$ with base weights and again uses the unweighted $q,q+1$ pool.
All final equations remain sums of squares of polynomial residuals.

\subsection{A parity argument replaces the removed factorial certificate}

For the moderate-coefficient source set
\[
 t=128+256u,\qquad K=t^4,
\]
and retain the request $F=t^K$. The smaller-source alternative sets
\[
 t=16+32u,\qquad K=2^{20}t.
\]
All parameters here range over all integers. In both cases $t$ is a nonzero
multiple of four. In the first case $K$ is positive and even; in the second,
the signed-power theorem forces $K\ge0$ and thus $t>0$. In either case the
first power gives a positive even $F$ with $F>|t|$.

The retained CRT row reads
\[
 Bv=C+F+1,\qquad C=2c.
\]
Its right-hand side is odd. Consequently $B$ is odd and nonzero, independently
of its eventual binomial interpretation. Also $H=2+4S$ for an integer $S$,
so $|BH|\ge2$. Its exact power gives
\[
 x\ge0,\qquad q=(BH)^x.
\]
Hence $q\ne0,-1$, and all five bases are nonzero. This proves that every
one of the old power requests is recovered by the new pools. It also justifies
the unweighted two-base variants, because the nonzero consecutive integers
$q$ and $q+1$ have the same sign.

\subsection{The endpoint removes the remaining negative horizon}

For the moderate-coefficient source, the sign of $t$ has not yet been
recovered. The retained beta endpoint rows give
\[
 y=F(2v+1),\qquad g=1+(t+1)y,\qquad x=1+g u_t.
\]
Since $|t|\ge128$ and $|y|\ge F$,
\[
 |g|\ge |t+1|F-1\ge127F-1>1.
\]
Thus $x=0$ would force $g u_t=-1$, a contradiction. It follows that $x\ge1$
and $|q|\ge2$. The literal retained request is $q^t$, with nominal index
$t$; its signed exactness forces $t\ge0$. Therefore $t\ge128$.

This argument uses neither a missing factorial row nor a false conclusion
from the binomial high quotient. In particular, setting $q=1$ in the digit
row does not identify $B$ alone: the high quotient remains. The endpoint
factorization is the required noncircular exclusion of $x=0$.

\subsection{Completeness and literal polynomial projections}

The positive branches of both affine progressions are cofinal and divisible
by four. Package 173's completeness construction can choose its horizon in
either progression. It needs precisely these properties and
$K\ge2^{20}t$, rather than the earlier square representation. The CRT anchor
$C$ can still be chosen positive, even, and larger than every other nominal
index. All subsequent cap and arithmetic witnesses can be completed.

Conversely the preceding sign deductions restore the hypotheses of the
retained arithmetic decoding theorem. The same label-separation and Horn-flow
arguments yield a finite C222 refutation, and the identified inherited
C222/ZF/ZFC calibration gives a finite ZFC refutation. Thus the new sources
have a zero exactly when the fixed ZFC presentation has a finite refutation.

Every subsequent projection uses a residual
$\epsilon v+f=0$ with $\epsilon\in\{1,-1\}$ and $v$ absent from $f$.
The supplied definition lists are acyclic. Substitution therefore gives a
bijection on integer zero sets, with the integral polynomial reconstruction
$v=-\epsilon f$. Registers in the final circuits are abbreviations and do
not add quantified variables.

The retained profiles are as follows. The last column bounds literal circuit
constants, not fully expanded polynomial coefficients.
\begin{center}
\begin{tabular}{rrrrr}
\hline
Variables & Exact degree & Residuals & Operations & Literal constant bound\\
\hline
131&12&70&2475&1099511627780\\
132&12&71&2484&256\\
117&16&56&2439&256\\
108&32&47&2434&256\\
97&128&36&2418&256\\
91&256&30&2400&256\\
87&32768&26&2388&256\\
86&2097780&25&2385&256\\
\hline
\end{tabular}
\end{center}

The immediate Oriented Powers comparisons include $137$ variables at degree
$12$, $121$ at degree $16$, $111$ at degree $32$, $99$ at degree $128$, and
$87$ at degree $8388612$. These are matched construction improvements, not
claims of global minimality. Exact degrees are certified by nonzero leading
homogeneous specializations; propagated coefficient-norm estimates provide
separate explicit upper bounds on expanded coefficient height.

The source checker independently binds all sixteen actual requests, fourteen
canonical Pell pairs, their shifted multiplication wires and norms, the
complete orientation row, and every squared frame atom. The circuit checker
verifies local polynomial expansions, acyclic unit pivots, source-row
coverage, degree certificates, and a nonzero derivative witness for each
quantified variable. These finite checks authenticate the construction;
they do not replace its unbounded arithmetic proof or the inherited logical
calibration.


\section{Why the consistency calibration is an ordinary ZFC proof}
\label{sec:calibration}

The direct and compact branches have different immediate calibrations.
The direct branch encodes the numerical proof calculus C222. It does
not simulate K215, instantiate Jones's universal polynomial, or use
Sun's integer-domain reduction. The compact branch deliberately retains
its existing native/K215 coefficient construction. The new argument
there is an exact soundness restoration to P66 and a fresh-witness
completeness construction; it does not replace the inherited coefficient
calibration by mere numerical inequalities.

\subsection{Direct branch}

The square-shell arithmetic constructors are
\begin{align*}
 I(a,b)&=2((a+b)^2+a),& M(a,b)&=I(a,b)+1,\\
 A(a,b)&=2((a+b)^2+3a+4b+3),&
 E(a,b)&=A(a,b)+1.
\end{align*}
They are the actual constructors in the arithmetic source.
Their constructor-preserving recoding to the original numerical C222
syntax fixes truth and falsity. One must not identify the two sets of
numeric constructors literally or assume that recoding preserves
arithmetic on their labels.

The inherited endpoints used after that recoding are
\begin{align*}
 &\texttt{KEndpoint.c222\_iff\_zf\_refutation},\\
 &\texttt{ConstructibleChoice.zf\_refutation\_iff\_zfc\_refutation}.
\end{align*}
The pinned source files are included for inspection and cited
directly~\cite{kendpoint,choice}. They contain the named theorem
compositions. Their full transitive import closure is not rebuilt by
this submission.

The lower C222 direction uses first-order completeness, named formulas
and a seed-to-model construction; the upper direction uses an
interpretation in a constructible model. The ZF/ZFC transport uses the
constructible submodel together with soundness and completeness. These
are ordinary constructions available in ZFC. The internally well-founded
part of the lower model is an \emph{internal} notion; no assumption of an
externally well-founded or transitive set model is added.

The arithmetic reductions are uniform over finite syntax, integers and
finite lists. Their inductions, finite CRT choices, Pell recurrences,
finite polynomial identities and finite group arguments are formalizable
in ZFC. Composing them with the identified logical endpoints gives the
internal equivalence in Theorem~\ref{thm:main}, without assuming consistency.
The entire inherited calibration is not described as a newly implemented
primitive-recursive proof translator: its supplied proof includes
semantic completeness and model constructions.

\subsection{Compact branch}

The four fixed coefficient ports are bound to one complete Start-only4127
native construction. Their full bytes, a finite reconstruction recipe
and a reconstructing program are supplied. Their total raw size is
97,046 bytes, and the largest fixed integer has 194,440 bits.
The new residue-ordinate proof reconstructs all missing positive P66
witnesses from a new zero. Conversely a refutation supplies the canonical
outer P66 witnesses, after which new private strong and auxiliary Pell
witnesses are chosen. This distinction matters: arbitrary old auxiliary
witnesses need not satisfy the new residue condition.

No coefficient port from another encoding is substituted.
The inequalities $\beta\ge31$, $0<M_C<\beta<M_F<2\beta$ are useful in
the proof, but they do not alone give a consistency calibration.
That calibration belongs to the whole named, authenticated tuple.

\subsection{Proof boundary}

The new unbounded arguments have received separately assigned
AI-assisted mathematical reviews. The checks supplied with the archive
authenticate finite algebra and exact source binding. They do not replace
the Pell classification theorem, the native window theorem, or the
C222/ZF/ZFC calibration by testing examples.

No new complete Lean or Rocq proof term or encoded object-language ZFC
derivation is supplied. This is an ordinary proof with exact computational
certificates and explicit inherited dependencies. It should be integrated
with that status rather than described as newly kernel-checked.

\section{Verification, reproducibility, and integration}
\label{sec:verify}

\subsection{What the programs check}

The quartic checker is independent of its emitter. It checks the exact
source bytes, the protected-row edits, commutation with the inherited
projection and optimization chain, the affine substitution, the two
invertible row operations, every coefficient of the collected sums of
squares, degree four and occurrence of every variable.
The signed horizon is a semantic change; the checker does not
mislabel it as a pointwise inverse to the old square horizon.

The compact checker reconstructs the literal graphs and their fixed
coefficient binding. It verifies all 56 terms of each leading form and
the 323-term main-norm cancellation. The integer-domain checker
authenticates the eleven positivity substitutions, the three free
integer coordinates, the 145-gate count, and the nonzero degree-288
leader. All four coefficient payloads are independently regenerated
from the complete fixed recipe and compared byte for byte.

The pooled-source checker binds sixteen actual power requests and fourteen
canonical Pell pairs, including every shifted multiplication wire,
norm, orientation term and required squared frame atom. The circuit
checker verifies all unit pivots and source-row partitions, local
polynomial expansions, degree witnesses, coefficient bounds and actual
dependence on each input. Deliberate corruptions of relevant artifact
fields are rejected. A fresh reconstruction reproduces its mathematical
JSON files exactly.

The release commands run the substantive checks both normally and under
\texttt{python -O}. Verification does not search for a zero and cannot
decide consistency. Degree statements use exact leading forms or
nonzero finite-field specializations of the highest homogeneous part,
not floating-point sampling.

\subsection{Files and commands}

The root directory contains this article, a root verifier, integration
notes, a SHA-256 manifest and three mathematical subpackages:
\texttt{quartic/}, \texttt{compact/} and \texttt{direct/}.
Each contains the relevant proof, literal equations and replay tools.
The final mathematical definitions are the exact integer residual lists,
straight-line circuits, fixed coefficient bytes and their stated domains.
Named intermediate registers are polynomial abbreviations, not hidden
quantified variables.

From the extracted package:
\begin{verbatim}
python3 -B verify_all.py
python3 -O -B verify_all.py
\end{verbatim}
Python 3.10 or later and its standard library suffice. The report can be
rebuilt with an ordinary TeX Live installation:
\begin{verbatim}
cd report
pdflatex -interaction=nonstopmode -halt-on-error report.tex
pdflatex -interaction=nonstopmode -halt-on-error report.tex
\end{verbatim}
PDF metadata can vary with the TeX installation. Frozen-byte manifests
and mathematical reconstruction have deliberately separate roles.
The integration notes suggest an additive new package; existing
calibrations and historical equations should remain available.

\section{Further research questions}
\label{sec:further}

\begin{enumerate}[leftmargin=*]
\item \textbf{A second residue in the same ordinate.}
Can an auxiliary Pell ordinate authenticate another quotient already
present in the compact formula without increasing the number of free
witnesses or losing the common-sign unit argument? The present proof
depends on $M\mid v$, $M\mid S^2$ and $R<M$; a second congruence would
have to be compatible with the same cofinal witness construction.

\item \textbf{Smaller carrier modulus.}
The $q^4$ carrier follows from the proven bound $R<q^4-2$.
A smaller modulus would need a sharper bound on the recovered
index difference, rather than an unsupported assertion that congruent
indices are equal.

\item \textbf{Sharper constructor growth budgets.}
The direct quartics use the uniform depth constant $\ell=2^{20}$.
For a nonnegative-coefficient polynomial $A(t)$ of degree at most $e$,
the finite inequality $A(t_0)<t_0^e$ implies
$A(t)<t^e$ for all $t\ge t_0$, by monotonicity of each term of
$A(t)/t^e$. Applying this to the actual constructor DAG can improve
its coarse depth budget. Extending a final horizon below eight
still requires rechecking every CRT and digit estimate; no such
additional equation is asserted here.

\item \textbf{Reduce quartic height after the sign change.}
The balanced $Q_{214}$ trades a variable for coefficient height.
Further integral coordinate or residual changes may reduce its
655,360 height without restoring a horizon variable. Any claimed
improvement should be measured on the fully collected polynomial.

\item \textbf{Simplify the logical source.}
The direct arithmetic cost is still driven by the large guarded C222
schema and its constructor tree. A different equiconsistent numerical
calculus might improve more than local certificate edits. Such a route
must supply its own uniform calibration, including the converse direction.

\item \textbf{Close the formalization boundary.}
The new sign recovery, residue-ordinate extraction, protected-row
commutation and exact circuit certificates are natural formalization
targets. Their finite algebra does not by itself close the inherited
unbounded arithmetic and logical dependency chain.
\end{enumerate}

\appendix
\section{Complete fixed coefficient recipe}\label{sec:ordinate-constants}
The following finite recipe is retained from the whole Start-only4127 construction~\cite{proveit174}. Its fixed indices do not introduce quantified variables into any equation. The accompanying four unsigned big-endian payloads have a combined size of 97,046 bytes.\par
\begingroup
\providecommand{\Cfl}[1]{\left\lfloor#1\right\rfloor}
\providecommand{\Crg}[2]{\{#1,\ldots,#2\}}
\textbf{The four fixed integers.}
Every index, set, sum, product and floor below defines a constant, not
another unknown. Put
\[
\begin{gathered}
\rho=2^{47},\quad B=\rho^{4127},\quad\beta=B-1,\\
G_m(x)=\sum_{j=0}^{m-1}x^j,\quad
\Lambda=\rho^{2042}+\rho^{3074}.
\end{gathered}
\]
\begin{align*}
M_C={}&\beta+\rho-(\rho+3)\rho^{990}\\
&-G_{943}(\rho)-\Psi(\rho,1),\\
M_F={}&B+3+(2^{44}-2)\Lambda\\
&+\rho^{973}G_9(\rho)+\rho^{1012}G_8(\rho)
  +\rho^{1031}\\
&+2^{29}\rho^{4124}-\rho^{4106},\\
K_0={}&\rho^{20}+\rho^{30}+\rho^{1031}+\rho^{4137}\\
&+\Lambda\Phi(1/\rho),
\end{align*}
where
\[
\begin{split}
\Phi(x)={}&\sum_{i=0}^{942}[1+2^{10}(1+w_i+2^8n_i\\
&+2^{17}p_i+2^{25}s_i)]x^i+2^{10}\Psi(x,2),\\
\Psi(x,t)={}&x^{992}(x^{10}+t^{17})G_8(tx)\\
&+t^8x^{943}(x^{30}+t^{17})G_9(tx)\\
&+t^{34}(2-t)x^{1031}.
\end{split}
\]
The rational intermediate in $K_0$ cancels exactly; all four results are positive integers.

\textbf{The 943 indexed fixed tiles.}
Use the actions $u_j,v_j,d_j$ decoded below. For $d=0,1$, let
\[
\begin{gathered}
h(q,b)=2q+b+2,\quad
\mathcal S_d=\{u_j:d_j=d,\ 0\le j<430\},\\
r_{d,q}=205-199d+|\mathcal S_d\cap\Crg{0}{q-1}|.
\end{gathered}
\]
An interval ending at $-1$ is empty. The table defines $t_i$ for
$0\le i\le942$; all index ranges are fixed:
\[
\setlength{\arraycolsep}{1.4pt}\begin{array}{c|l}
\text{index }i&\text{tile }t_i=(w_i,n_i,p_i,s_i)\\\hline
0,1&(2,435,3,1),\ (5,0,0,0)\\
j+2\ (j=2q+b<430)&((1-d_j)r_{0,u_j},h(q,b),d_jr_{1,u_j},v_j)\\
434-2a+b\ (a,b=0,1)&(0,432a+b,0,432a+b)\\
\substack{424+2r_{d,q}+b\\q\in\mathcal S_d,\ d,b=0,1}&(dr_{d,q},b,(1-d)r_{d,q},h(q,b))\\
924+2a+b\ (a,b=0,1)&(0,432a+b,0,438)\\
928&(2,434,2,0)\\
929&(3,436,4,4)\\
930&(4,436,4,0)\\
497+t\ (t=434,435,436)&(1,438,1,t)\\
934+3a+b\ (a,b=0,1,2)&(a^2,437,b,437)
\end{array}
\]
\textbf{A closed decoder for the 430 actions.}
Put $D_b(A,k)=\Cfl{A/b^k}\bmod b$, where
$a\bmod b=a-b\Cfl{a/b}$. For $0\le q<215$, $b=0,1$, $j=2q+b$:
\[
\begin{gathered}
N=\mathcal A\bmod2^{636},\quad
T_0=\Cfl{\mathcal A/2^{636}}\bmod2^{1675},\\
E_0=\Cfl{\mathcal A/2^{2311}},\\
\nu_j=D_2(N,j),\quad \sigma_j=\sum_{k=0}^{j-1}\nu_k,\\
m_q=D_2(N,430+q),\quad \tau_q=\sum_{k=0}^{q-1}m_k,\\
F_q=14+m_q(D_{16}(E_0,\tau_q)-14),\\
u_j=\nu_j(1+\sigma_j)+(1-\nu_j)D_{216}(T_0,j-\sigma_j),\\
v_j=D_2(F_q,2b),\quad d_j=D_2(F_q,2b+1).
\end{gathered}
\]
Empty sums are zero. Here $(v_0,d_0,u_0)=(1,1,1)$; the constructor tile
uses this actual first target, including its state renaming.

\textbf{The integer $\mathcal A$.}
The alphabet, in order of digit values $0,\ldots,61$, is:\par
\noindent{\small\ttfamily 0123456789ABCDEFGHIJKLMNOPQRSTUVWXYZabcdefghijklmnopqrstuvwxyz}\par
Concatenate the following 436 characters as digits $e_0,\ldots,e_{435}$;
put $\mathcal A=\sum_{j=0}^{435}e_j62^{435-j}$.\par
\begingroup\small\ttfamily
BxEnhvy6NipzkWmROo0GsA3uK3twWLxscLq20UTumy5Cdu461LwD6g7RDNzF8Y\par
iYDFTpB8ehAypd7d6OsOL4APamnJeNXQypUezCmW9n4oXP6XfgmiWuAWZ8ZJ7U\par
OP7gu4kVdlGZqZHULQMar1dXNeJW4i7oNVI8usp8BlovjjZbwPCovZI328fRUq\par
4q92R81s5GjWQaQfGhkq9UYMK3rRMaBmsOIBkJqc2cj2VRIZ6mjzkBIOkW9SjZ\par
FzeA6J5u4j3sBL3SdgwXzM9mHvQivu8EJ1TBTpRhHDqZIbmYgwQ8VqTAFotl6F\par
We1YC554GKrVHsA8oUGKSA6F6k5KREBSD0Tgd7bFALY84LLbpw574BQyvdoHrl\par
0Oe35uwDf15DpT6VJWZBjSbHPpPAINYTte0eUxnTFTKBlGj8Am6sb1gLldh8WZ\par
8P\par
\endgroup

\endgroup


\clearpage
\section{Complete literal balanced quartic}\label{app:quartic}

Every $x_0,\ldots,x_{213}$ below ranges independently over $\Z$.
The equation is
\[
 \boxed{Q_{214}(x_0,\ldots,x_{213})=\sum_{j=0}^{152}R_j^2=0.}
\]
The list is generated directly from the verified integer coefficient file.
It contains 153 quadratic residuals; there are no unstated side
conditions or auxiliary quantifiers. The adjacent JSON also records
every collected coefficient of the one quartic.

% Generated by emit_quartic_appendix.py; do not hand-edit.
% Source SHA-256: e85fa439ed6a93835c17e326460224f9f04ea7a54bfc7175a25678fa5a637318
\begingroup\small
\setlength{\jot}{1.5pt}
\begin{align*}
R_{0} &= 8 -9x_{3} -16x_{3}x_{184} +x_{6} +16x_{184}\\
R_{1} &= -x_{1} -x_{3}x_{4} +x_{5}x_{6}\\
R_{2} &= -1 +x_{0} -x_{5}x_{6}\\
R_{3} &= -2x_{1}x_{212} -x_{2}x_{17} +x_{3} +2x_{15}\\
R_{4} &= -2 -x_{7} -x_{7}^{2} -2x_{7}x_{8} -2x_{8}\\
&\quad -x_{8}^{2} +x_{18}\\
R_{5} &= -2 -x_{7} -x_{7}^{2} -8x_{7}x_{18} -4x_{7}x_{19}\\
&\quad -8x_{18} -16x_{18}^{2} -16x_{18}x_{19} -4x_{19} -4x_{19}^{2}\\
&\quad +x_{20}\\
R_{6} &= 1 +2x_{20} -4x_{20}^{2} +x_{21}\\
R_{7} &= -2 -x_{8}^{2} -2x_{8}x_{9} -x_{9}^{2} +x_{22}\\
R_{8} &= -1 +2x_{7} -x_{7}^{2} -2x_{7}x_{9} +2x_{9}\\
&\quad -x_{9}^{2} +x_{23}\\
R_{9} &= -7 +33x_{7} -35x_{7}^{2} -96x_{7}x_{9} -24x_{7}x_{18}\\
&\quad -4x_{7}x_{22} -24x_{7}x_{23} +16x_{8} -2x_{8}^{2} -8x_{8}x_{22}\\
&\quad +32x_{9} -64x_{9}^{2} -32x_{9}x_{18} -32x_{9}x_{23} +4x_{18}\\
&\quad -4x_{18}^{2} -8x_{18}x_{23} +12x_{22} -4x_{22}^{2} +8x_{23}\\
&\quad -4x_{23}^{2} +x_{24}\\
R_{10} &= -16 +14x_{7} -2x_{7}^{2} -8x_{7}x_{8} -8x_{7}x_{22}\\
&\quad +8x_{8} +8x_{22} -4x_{24}^{2} +x_{25} -2x_{42}\\
R_{11} &= -6 +16x_{9} -16x_{9}^{2} -16x_{9}x_{22} -16x_{9}x_{23}\\
&\quad +8x_{22} -4x_{22}^{2} -8x_{22}x_{23} +8x_{23} -4x_{23}^{2}\\
&\quad +x_{26}\\
R_{12} &= -5 +24x_{9} -16x_{9}^{2} -16x_{9}x_{18} -32x_{9}x_{23}\\
&\quad -16x_{9}x_{26} +10x_{18} -4x_{18}^{2} -16x_{18}x_{23} -8x_{18}x_{26}\\
&\quad +24x_{23} -16x_{23}^{2} -16x_{23}x_{26} +12x_{26} -4x_{26}^{2}\\
&\quad +x_{27}\\
R_{13} &= -12 -29x_{7} -37x_{7}^{2} -24x_{7}x_{8} -112x_{7}x_{9}\\
&\quad -4x_{7}x_{22} -24x_{7}x_{23} -39x_{8} -37x_{8}^{2} -112x_{8}x_{9}\\
&\quad -24x_{8}x_{22} -4x_{8}x_{23} -80x_{9} -128x_{9}^{2} -32x_{9}x_{22}\\
&\quad -32x_{9}x_{23} -12x_{22} -4x_{22}^{2} -8x_{23} -4x_{23}^{2}\\
&\quad +x_{28}\\
R_{14} &= -18 -10x_{7} -2x_{7}^{2} -24x_{7}x_{8} -32x_{7}x_{9}\\
&\quad -8x_{7}x_{22} -72x_{8} -72x_{8}^{2} -192x_{8}x_{9} -48x_{8}x_{22}\\
&\quad -96x_{9} -128x_{9}^{2} -64x_{9}x_{22} -24x_{22} -8x_{22}^{2}\\
&\quad -4x_{28}^{2} +x_{29}\\
R_{15} &= -32 -42x_{7} +8x_{7}x_{9} +16x_{7}x_{22} -84x_{8}\\
&\quad +16x_{8}x_{9} +32x_{8}x_{22} +21x_{9} -3x_{9}^{2} -4x_{9}x_{18}\\
&\quad -8x_{9}x_{22} +21x_{18} -8x_{18}x_{22} +24x_{22} -4x_{22}^{2}\\
&\quad +x_{30} -x_{58}\\
R_{16} &= -1 -24x_{7} -64x_{7}^{2} -160x_{7}x_{9} -64x_{7}x_{22}\\
&\quad -32x_{7}x_{23} -32x_{7}x_{30} -26x_{9} -100x_{9}^{2} -80x_{9}x_{22}\\
&\quad -40x_{9}x_{23} -40x_{9}x_{30} -8x_{22} -16x_{22}^{2} -16x_{22}x_{23}\\
&\quad -16x_{22}x_{30} -6x_{23} -4x_{23}^{2} -8x_{23}x_{30} -4x_{30}\\
&\quad -4x_{30}^{2} +x_{31}\\
R_{17} &= 1 -4x_{18}^{2} +x_{19}\\
R_{18} &= -2 -x_{7} -x_{7}^{2} +8x_{7}x_{18} -4x_{7}x_{19}\\
&\quad +8x_{18} -16x_{18}^{2} +16x_{18}x_{19} -4x_{19} -4x_{19}^{2}\\
&\quad +x_{32}\\
R_{19} &= 1 +2x_{32} -4x_{32}^{2} +x_{33}\\
R_{20} &= -x_{8}^{2} -4x_{8}x_{33} -4x_{33}^{2} +x_{34}\\
R_{21} &= -8 -20x_{7} -16x_{7}^{2} -64x_{7}x_{8} +32x_{7}x_{18}\\
&\quad -40x_{8} -64x_{8}^{2} +64x_{8}x_{18} +18x_{18} -16x_{18}^{2}\\
&\quad +x_{35}\\
R_{22} &= 8x_{8} -62x_{8}^{2} -92x_{8}x_{9} -32x_{8}x_{22} -32x_{8}x_{35}\\
&\quad +6x_{9} -34x_{9}^{2} -24x_{9}x_{22} -24x_{9}x_{35} -4x_{22}^{2}\\
&\quad -8x_{22}x_{35} +4x_{35} -4x_{35}^{2} +x_{36}\\
R_{23} &= 8x_{7} -64x_{7}^{2} -96x_{7}x_{9} -32x_{7}x_{23} -32x_{7}x_{36}\\
&\quad +6x_{9} -36x_{9}^{2} -24x_{9}x_{23} -24x_{9}x_{36} +2x_{23}\\
&\quad -4x_{23}^{2} -8x_{23}x_{36} +4x_{36} -4x_{36}^{2} +x_{37}\\
R_{24} &= -x_{9}^{2} -2x_{9}x_{10} -x_{10}^{2} +x_{38}\\
R_{25} &= -32 +22x_{7} -3x_{7}^{2} -8x_{7}x_{8} -8x_{7}x_{22}\\
&\quad -8x_{7}x_{23} -4x_{7}x_{38} +16x_{8} -8x_{8}x_{23} -2x_{9}^{2}\\
&\quad -8x_{9}x_{38} +16x_{22} -8x_{22}x_{23} +22x_{23} -4x_{23}^{2}\\
&\quad +4x_{38} -4x_{38}^{2} +x_{39} -x_{42}\\
R_{26} &= 6 -16x_{7}^{2} -36x_{7}x_{9} -32x_{7}x_{23} -40x_{7}x_{38}\\
&\quad -16x_{7}x_{39} -22x_{9}^{2} -32x_{9}x_{23} -48x_{9}x_{38} -16x_{9}x_{39}\\
&\quad -2x_{23} -16x_{23}^{2} -32x_{23}x_{38} -16x_{23}x_{39} +4x_{38}\\
&\quad -24x_{38}^{2} -16x_{38}x_{39} -4x_{39}^{2} +x_{40}\\
R_{27} &= x_{9} -25x_{9}^{2} -40x_{9}x_{23} -20x_{9}x_{26} -16x_{23}^{2}\\
&\quad -16x_{23}x_{26} -4x_{26}^{2} +x_{41}\\
R_{28} &= -1 +8x_{8} -4x_{8}^{2} -8x_{8}x_{22} +8x_{22}\\
&\quad -4x_{22}^{2} +x_{42}\\
R_{29} &= -1 -16x_{8}^{2} -8x_{8}x_{9} -32x_{8}x_{22} -16x_{8}x_{42}\\
&\quad -x_{9}^{2} -8x_{9}x_{22} -4x_{9}x_{42} -16x_{22}^{2} -16x_{22}x_{42}\\
&\quad -4x_{42}^{2} +x_{43}\\
R_{30} &= 3 +x_{8} -2x_{8}x_{9} -4x_{8}x_{43} -3x_{9}^{2}\\
&\quad -8x_{9}x_{43} -x_{22} -4x_{43}^{2} +x_{44}\\
R_{31} &= -1 -4x_{44} -4x_{44}^{2} +x_{45}\\
R_{32} &= -1 -x_{9}^{2} -8x_{9}x_{44} -4x_{9}x_{45} -16x_{44}^{2}\\
&\quad -16x_{44}x_{45} -4x_{45}^{2} +x_{46}\\
R_{33} &= 2 +x_{8} -x_{8}^{2} -4x_{8}x_{9} -4x_{8}x_{46}\\
&\quad -4x_{9}^{2} -8x_{9}x_{46} -4x_{46}^{2} +x_{47}\\
R_{34} &= -4x_{41}^{2} -8x_{41}x_{47} -4x_{47}^{2} +x_{48}\\
R_{35} &= -9 +7x_{7} -x_{7}^{2} -8x_{7}x_{41} -4x_{7}x_{48}\\
&\quad +24x_{41} -16x_{41}^{2} -16x_{41}x_{48} +12x_{48} -4x_{48}^{2}\\
&\quad +x_{49}\\
R_{36} &= -1 -4x_{49} -4x_{49}^{2} +x_{50}\\
R_{37} &= -x_{9}^{2} -8x_{9}x_{49} -4x_{9}x_{50} -16x_{49}^{2} -16x_{49}x_{50}\\
&\quad -4x_{50}^{2} +x_{51}\\
R_{38} &= -x_{7} -x_{7}^{2} -4x_{7}x_{9} -4x_{7}x_{51} -4x_{9}\\
&\quad -4x_{9}^{2} -8x_{9}x_{51} -4x_{51} -4x_{51}^{2} +x_{52}\\
R_{39} &= -10 -12x_{52} -4x_{52}^{2} +x_{53}\\
R_{40} &= -19 +4x_{8} -4x_{8}^{2} -8x_{8}x_{22} -16x_{8}x_{52}\\
&\quad -8x_{8}x_{53} +4x_{22} -4x_{22}^{2} -16x_{22}x_{52} -8x_{22}x_{53}\\
&\quad -16x_{52} -24x_{52}^{2} -16x_{52}x_{53} +6x_{53} -4x_{53}^{2}\\
&\quad +x_{54}\\
R_{41} &= -10 +x_{8} -25x_{8}^{2} -40x_{8}x_{22} -20x_{8}x_{54}\\
&\quad -16x_{22}^{2} -16x_{22}x_{54} -12x_{47} -4x_{47}^{2} -4x_{54}^{2}\\
&\quad +x_{55}\\
R_{42} &= 21 +24x_{47} -8x_{47}^{2} -16x_{47}x_{55} -2x_{55}\\
&\quad -4x_{55}^{2} +x_{56}\\
R_{43} &= -5 +4x_{7} +8x_{8} +x_{9} -x_{9}^{2}\\
&\quad -4x_{9}x_{56} -2x_{18} -4x_{56}^{2} +x_{57} +12x_{58}\\
&\quad -4x_{58}^{2}\\
R_{44} &= -6x_{7} -16x_{7}^{2} -64x_{7}x_{8} +8x_{7}x_{9} +16x_{7}x_{18}\\
&\quad -12x_{8} -64x_{8}^{2} +16x_{8}x_{9} +32x_{8}x_{18} +3x_{9}\\
&\quad -x_{9}^{2} -4x_{9}x_{18} +3x_{18} -4x_{18}^{2} +x_{58}\\
R_{45} &= -12 +8x_{7} +16x_{8} -4x_{18} -4x_{57}^{2}\\
&\quad +24x_{58} -8x_{58}^{2} +x_{59}\\
R_{46} &= -1 -4x_{59} -4x_{59}^{2} +x_{60}\\
R_{47} &= -x_{1}x_{12} +x_{61}\\
R_{48} &= x_{11} -x_{12} -2x_{21}x_{61} -x_{62}\\
R_{49} &= x_{11} -x_{12} -2x_{25}x_{61} -x_{63}\\
R_{50} &= x_{11} -x_{12} -2x_{27}x_{61} -6x_{61} -x_{64}\\
R_{51} &= x_{11} -x_{12} -2x_{29}x_{61} -x_{65}\\
R_{52} &= x_{11} -x_{12} -2x_{31}x_{61} +2x_{61} -x_{66}\\
R_{53} &= x_{11} -x_{12} -4x_{33}x_{61} -2x_{34}x_{61} -x_{67}\\
R_{54} &= x_{11} -x_{12} -2x_{37}x_{61} -x_{68}\\
R_{55} &= 2x_{27}x_{61} -2x_{40}x_{61} -6x_{61} +x_{64} -x_{69}\\
R_{56} &= -x_{7}x_{8} +x_{7}x_{9} +x_{8} +x_{8}x_{9} -x_{9}\\
&\quad -x_{9}^{2} +x_{70}\\
R_{57} &= 1 +2x_{10} +x_{10}^{2} -x_{70}^{2} +x_{71}\\
R_{58} &= -x_{11}^{2} +2x_{11}x_{12} +2x_{11}x_{73} -x_{12}^{2} -2x_{12}x_{73}\\
&\quad -x_{71}^{2} +x_{72} -x_{73}^{2}\\
R_{59} &= -4x_{59}x_{61} -2x_{60}x_{61} +x_{73}\\
R_{60} &= -x_{8}x_{61} -x_{12} +x_{74}x_{75} +x_{76}\\
R_{61} &= -1 +4x_{1} +4x_{1}x_{7} +8x_{1}x_{8} -2x_{1}x_{18}\\
&\quad +x_{74}\\
R_{62} &= -x_{7}x_{207} -x_{11} +x_{75} +x_{207}\\
R_{63} &= x_{8}x_{207} -2x_{18}x_{61} +2x_{31}x_{61} +x_{66} +x_{77}\\
R_{64} &= -x_{11}x_{62} +x_{12}x_{62} +x_{78}\\
R_{65} &= -x_{63}x_{78} +x_{79}\\
R_{66} &= -x_{64}x_{79} +x_{80}\\
R_{67} &= -x_{65}x_{80} +x_{81}\\
R_{68} &= -x_{66}x_{81} +x_{82}\\
R_{69} &= -x_{67}x_{82} +x_{83}\\
R_{70} &= -x_{68}x_{83} +x_{84}\\
R_{71} &= -x_{69}x_{84} +x_{85}\\
R_{72} &= -x_{76}x_{85} +x_{86}\\
R_{73} &= -x_{77}x_{86} +x_{87}\\
R_{74} &= -x_{87}^{2} +x_{88}\\
R_{75} &= -x_{3} -2x_{3}x_{15} +2x_{15} +x_{89}\\
R_{76} &= -x_{13}x_{89} +x_{90}\\
R_{77} &= -2x_{0}^{2} +2x_{0}x_{11} +2x_{0}x_{12} +2x_{0}x_{90} +2x_{0}x_{91}\\
&\quad -x_{11}^{2} -2x_{11}x_{90} -x_{12}^{2} -2x_{12}x_{91} +x_{16}x_{93}\\
&\quad -x_{72}x_{88} -x_{90}^{2} -x_{91}^{2}\\
R_{78} &= -x_{3}x_{14} +x_{14} -x_{14}x_{89} +x_{91}\\
R_{79} &= -x_{94}x_{97} +x_{97} +x_{99}\\
R_{80} &= -x_{2}x_{120} +x_{95}x_{199} -x_{98} -x_{99}x_{120}\\
R_{81} &= x_{100}x_{199} -x_{101}x_{123} -x_{102}\\
R_{82} &= x_{103}x_{199} -x_{104}x_{123} -x_{105}\\
R_{83} &= x_{106}x_{199} -x_{107}x_{123} -x_{108}\\
R_{84} &= x_{109}x_{199} -x_{110}x_{123} -x_{111}\\
R_{85} &= x_{112}x_{200} -x_{114}x_{123} -x_{115}\\
R_{86} &= x_{116}x_{200} -x_{117}x_{123} -x_{118}\\
R_{87} &= -x_{16}x_{120} +x_{123}\\
R_{88} &= x_{119} -x_{121}x_{123} -x_{122}\\
R_{89} &= -x_{123}x_{125} +x_{124} -x_{126}\\
R_{90} &= -x_{123}x_{128} +x_{127}x_{199} -x_{129}\\
R_{91} &= x_{130} -2x_{130}x_{132} +x_{133}\\
R_{92} &= -x_{130}x_{133} -x_{133} +x_{134}\\
R_{93} &= -1 +x_{131}^{2} +x_{133}^{2} -x_{134}^{2}\\
R_{94} &= 1 -x_{131}x_{136} -x_{135} +x_{137}\\
R_{95} &= -1 +2x_{135} -x_{135}^{2} +2x_{135}x_{137} +x_{136}^{2}\\
R_{96} &= 65 -16x_{131} -32x_{131}x_{184} +x_{138} +256x_{184}\\
&\quad +256x_{184}^{2}\\
R_{97} &= x_{1} +x_{138}x_{139} -x_{180} -8x_{181} -16x_{181}x_{184}\\
R_{98} &= 1 +x_{2}^{2} -2x_{2}x_{131} +x_{140}\\
R_{99} &= -x_{2}x_{172} +x_{16} +x_{140}x_{141} -x_{171}\\
R_{100} &= -1 +2x_{131}x_{183} -x_{142} -x_{183}^{2}\\
R_{101} &= -1 +x_{94} +x_{142}x_{143} -x_{171} -x_{172}x_{183}\\
R_{102} &= -x_{131}x_{145} -x_{144} +x_{146}\\
R_{103} &= -1 -x_{144}^{2} +2x_{144}x_{146} +x_{145}^{2}\\
R_{104} &= 1 -x_{131}x_{148} -x_{147} +x_{149}\\
R_{105} &= -1 +2x_{147} -x_{147}^{2} +2x_{147}x_{149} +x_{148}^{2}\\
R_{106} &= -x_{131}x_{151} -x_{150} +x_{152}\\
R_{107} &= -1 -x_{150}^{2} +2x_{150}x_{152} +x_{151}^{2}\\
R_{108} &= -x_{131}x_{154} -x_{153} +x_{155}\\
R_{109} &= -1 -x_{153}^{2} +2x_{153}x_{155} +x_{154}^{2}\\
R_{110} &= -x_{131}x_{157} -x_{156} +x_{158}\\
R_{111} &= -1 -x_{156}^{2} +2x_{156}x_{158} +x_{157}^{2}\\
R_{112} &= -x_{131}x_{160} -x_{159} +x_{161}\\
R_{113} &= -1 -x_{159}^{2} +2x_{159}x_{161} +x_{160}^{2}\\
R_{114} &= 1 -x_{131}x_{163} -x_{162} +x_{164}\\
R_{115} &= -1 +2x_{162} -x_{162}^{2} +2x_{162}x_{164} +x_{163}^{2}\\
R_{116} &= -x_{131}x_{166} -x_{165} +x_{167}\\
R_{117} &= -1 -x_{165}^{2} +2x_{165}x_{167} +x_{166}^{2}\\
R_{118} &= -x_{131}x_{169} -x_{168} +x_{170}\\
R_{119} &= -1 -x_{168}^{2} +2x_{168}x_{170} +x_{169}^{2}\\
R_{120} &= 1 -x_{131}x_{172} -x_{171} +x_{173}\\
R_{121} &= -1 +2x_{171} -x_{171}^{2} +2x_{171}x_{173} +x_{172}^{2}\\
R_{122} &= -x_{131}x_{175} -x_{174} +x_{176}\\
R_{123} &= -1 -x_{174}^{2} +2x_{174}x_{176} +x_{175}^{2}\\
R_{124} &= -x_{131}x_{178} -x_{177} +x_{179}\\
R_{125} &= -1 -x_{177}^{2} +2x_{177}x_{179} +x_{178}^{2}\\
R_{126} &= 1 -x_{131}x_{181} -x_{180} +x_{182}\\
R_{127} &= -1 +2x_{180} -x_{180}^{2} +2x_{180}x_{182} +x_{181}^{2}\\
R_{128} &= -21 -x_{0}^{2} -16x_{1} -2x_{1}^{2} -32x_{1}x_{184}\\
&\quad -x_{2}^{2} -x_{3}^{2} +2x_{7} -x_{7}^{2} -x_{8}^{2}\\
&\quad -x_{9}^{2} -x_{10}^{2} -x_{11}^{2} -x_{12}^{2} -x_{13}^{2}\\
&\quad -x_{14}^{2} -4x_{15}^{2} -x_{16}^{2} -x_{92}^{2} +2x_{94}\\
&\quad -2x_{94}^{2} -x_{95}^{2} -x_{96}^{2} -x_{100}^{2} -x_{103}^{2}\\
&\quad -x_{106}^{2} -x_{109}^{2} -x_{112}^{2} -x_{113}^{2} -x_{116}^{2}\\
&\quad -x_{119}^{2} -x_{120}^{2} -x_{124}^{2} -x_{127}^{2} +x_{130}\\
&\quad -x_{183}^{2} -32x_{184} -2x_{213}\\
R_{129} &= 2x_{2} -4x_{2}x_{7} -4x_{2}x_{8} -4x_{2}x_{9} -4x_{2}x_{10}\\
&\quad -4x_{2}x_{11} -4x_{2}x_{12} -4x_{2}x_{13} -4x_{2}x_{14} -8x_{2}x_{15}\\
&\quad +x_{183}\\
R_{130} &= -x_{131}x_{186} +x_{195}\\
R_{131} &= -1 +x_{185}^{2} +x_{186}^{2} -x_{195}^{2}\\
R_{132} &= -x_{188}x_{189} +x_{196}\\
R_{133} &= -1 +x_{187}^{2} +x_{188}^{2} -x_{196}^{2}\\
R_{134} &= -x_{178}^{2} +x_{197}\\
R_{135} &= x_{186} -2x_{191}x_{197} -x_{197}\\
R_{136} &= -1 -4x_{186} -8x_{186}x_{190} +x_{189}\\
R_{137} &= -x_{131} -x_{185}x_{192} +x_{189}\\
R_{138} &= -x_{179} -x_{185}x_{193} +x_{187}\\
R_{139} &= -2x_{15} -4x_{178}x_{194} +x_{188}\\
R_{140} &= 9 +x_{0}^{2} -2x_{0}x_{172} +16x_{1} +x_{1}^{2}\\
&\quad -2x_{1}x_{145} +32x_{1}x_{184} +x_{3}^{2} -2x_{3}x_{163} -2x_{7}\\
&\quad +x_{7}^{2} -2x_{7}x_{148} +x_{8}^{2} -2x_{8}x_{151} +x_{9}^{2}\\
&\quad -2x_{9}x_{154} +x_{10}^{2} -2x_{10}x_{157} +x_{11}^{2} -2x_{11}x_{160}\\
&\quad +x_{12}^{2} -2x_{12}x_{166} +x_{13}^{2} -2x_{13}x_{169} +x_{14}^{2}\\
&\quad -2x_{14}x_{175} -4x_{15}x_{178} +x_{92}^{2} -2x_{92}x_{181} -x_{131}x_{198}\\
&\quad -16x_{136} -32x_{136}x_{184} +x_{137}^{2} -16x_{145} -32x_{145}x_{184}\\
&\quad -2x_{146} +x_{146}^{2} +2x_{148} +x_{149}^{2} -2x_{152}\\
&\quad +x_{152}^{2} -2x_{155} +x_{155}^{2} -2x_{158} +x_{158}^{2}\\
&\quad -2x_{161} +x_{161}^{2} +x_{164}^{2} -2x_{167} +x_{167}^{2}\\
&\quad -2x_{170} +x_{170}^{2} +x_{173}^{2} -2x_{176} +x_{176}^{2}\\
&\quad +2x_{178}^{2} -2x_{179} +x_{179}^{2} +x_{182}^{2} +x_{198}\\
&\quad -x_{201}^{2} -x_{202}^{2} -x_{203} -x_{203}^{2} +x_{213}\\
R_{141} &= -x_{96}x_{199} +x_{120}\\
R_{142} &= -x_{113}x_{200} +x_{120}\\
R_{143} &= -1 -4x_{15}^{2} -x_{136}^{2} -x_{145}^{2} -x_{148}^{2}\\
&\quad -x_{151}^{2} -x_{154}^{2} -x_{157}^{2} -x_{160}^{2} -x_{163}^{2}\\
&\quad -x_{166}^{2} -x_{169}^{2} -x_{172}^{2} -x_{175}^{2} +x_{178}^{2}\\
&\quad -x_{181}^{2} -x_{201}^{2} -x_{202}^{2} -x_{203} -x_{203}^{2}\\
R_{144} &= -1 -4x_{98}^{2} -4x_{102}^{2} -4x_{105}^{2} -4x_{108}^{2}\\
&\quad -4x_{111}^{2} -4x_{115}^{2} -4x_{118}^{2} +x_{120}^{2} -4x_{122}^{2}\\
&\quad -4x_{126}^{2} -4x_{129}^{2} -x_{204}^{2} -x_{205}^{2} -x_{206}\\
&\quad -x_{206}^{2}\\
R_{145} &= -x_{1}x_{11} +x_{207}\\
R_{146} &= x_{94}x_{209} -2x_{95}x_{145} -2x_{100}x_{148} -2x_{103}x_{151} -2x_{106}x_{154}\\
&\quad -2x_{109}x_{157} -2x_{112}x_{160} -2x_{116}x_{166} -2x_{119}x_{169} -2x_{124}x_{175}\\
&\quad -2x_{127}x_{178} -2x_{131}x_{209} +2x_{145}x_{146} +2x_{148} +2x_{148}x_{149}\\
&\quad +2x_{151}x_{152} +2x_{154}x_{155} +2x_{157}x_{158} +2x_{160}x_{161} +2x_{166}x_{167}\\
&\quad +2x_{169}x_{170} +2x_{175}x_{176} +2x_{178}x_{179} +x_{208}\\
R_{147} &= -10 -x_{94}x_{208} +x_{95}^{2} -2x_{95}x_{144} +x_{100}^{2}\\
&\quad -2x_{100}x_{147} +x_{103}^{2} -2x_{103}x_{150} +x_{106}^{2} -2x_{106}x_{153}\\
&\quad +x_{109}^{2} -2x_{109}x_{156} +x_{112}^{2} -2x_{112}x_{159} +x_{116}^{2}\\
&\quad -2x_{116}x_{165} +x_{119}^{2} -2x_{119}x_{168} +x_{124}^{2} -2x_{124}x_{174}\\
&\quad +x_{127}^{2} -2x_{127}x_{177} +2x_{144}x_{146} +2x_{147} +2x_{147}x_{149}\\
&\quad +2x_{150}x_{152} +2x_{153}x_{155} +2x_{156}x_{158} +2x_{159}x_{161} +2x_{165}x_{167}\\
&\quad +2x_{168}x_{170} +2x_{174}x_{176} +2x_{177}x_{179} +x_{209}\\
R_{148} &= x_{94}x_{211} -2x_{96}x_{136} -2x_{113}x_{163} -2x_{120}x_{172} -2x_{131}x_{211}\\
&\quad +2x_{136} +2x_{136}x_{137} +2x_{163} +2x_{163}x_{164} +2x_{172}\\
&\quad +2x_{172}x_{173} +x_{210} -x_{211}\\
R_{149} &= -2 -x_{94}x_{210} +x_{96}^{2} -2x_{96}x_{135} +x_{113}^{2}\\
&\quad -2x_{113}x_{162} +x_{120}^{2} -2x_{120}x_{171} +2x_{135} +2x_{135}x_{137}\\
&\quad +2x_{162} +2x_{162}x_{164} +x_{171}^{2} -x_{172}^{2} +x_{210}\\
&\quad +x_{211}\\
R_{150} &= -1 -x_{1} -2x_{1}x_{212} +x_{3}\\
R_{151} &= -128 -512x_{184} -512x_{184}^{2} +x_{213}\\
R_{152} &= x_{92} -512x_{213}^{2}
\end{align*}
\endgroup


\clearpage
\section{Source identities and comparison notes}

The exact repository source identities, original Git blob identifiers and
SHA-256 digests are retained in each subpackage's provenance ledger.
The earlier Oriented Powers archive has SHA-256
\begin{center}
\small\texttt{b5f564d8afc22cd0d17fd831f965baf9fa3423155229c7caa9b1084cddcbbaa2}.
\end{center}
The final archive's manifest authenticates its delivered contents.
Inherited notices identify prior authorship and source scope.
The new manuscript and code are AI-assisted work prepared for ProveIt;
separate AI reviews are not described as external human peer review.

\begin{thebibliography}{99}
\raggedright
\bibitem{proveit173}
ProveIt contributors, \emph{A reused CRT index: an explicit Diophantine
consistency equation}, package 173, 25 October 2026.
\url{https://github.com/VladimirReshetnikov/ProveIt/tree/02a78a08603cbae248c19c23a75907c6c6bd6443/Computability/DiophantineConsistency/con-zfc-diophantine-forms/173-crt-index-anchor-215}.
\bibitem{proveit174}
ProveIt contributors, \emph{Start-marker compositions}, package 174,
25 October 2026. Complete Start-only4127 tuple and P66 calibration.
\url{https://github.com/VladimirReshetnikov/ProveIt/tree/02a78a08603cbae248c19c23a75907c6c6bd6443/Computability/DiophantineConsistency/con-zfc-diophantine-forms/174-start-marker-compositions}.
\bibitem{proveit170}
ProveIt contributors, \emph{Quotient carriers}, package 170,
25 October 2026.
\url{https://github.com/VladimirReshetnikov/ProveIt/tree/02a78a08603cbae248c19c23a75907c6c6bd6443/Computability/DiophantineConsistency/con-zfc-diophantine-forms/170-quotient-carriers}.
\bibitem{oriented}
\emph{Oriented power pooling and polynomial root centering},
AI-assisted research prepared for Vladimir Reshetnikov and ProveIt,
25 October 2026. Prior delivered archive
\texttt{ConZFC\_Oriented\_Powers\_2026-10-25.zip}; exact identity above.
The local inherited theorem and source pins are in \texttt{direct/}.
\bibitem{kendpoint}
ProveIt contributors, \texttt{KEndpoint.lean},
\texttt{c222\_iff\_zf\_refutation}.
\url{https://github.com/VladimirReshetnikov/ProveIt/blob/02a78a08603cbae248c19c23a75907c6c6bd6443/Computability/BusyBeaver/Lean/BusyBeaver/ZFC/KEndpoint.lean}.
\bibitem{choice}
ProveIt contributors, \texttt{ConstructibleChoice.lean},
\texttt{zf\_refutation\_iff\_zfc\_refutation}.
\url{https://github.com/VladimirReshetnikov/ProveIt/blob/02a78a08603cbae248c19c23a75907c6c6bd6443/Computability/BusyBeaver/Lean/BusyBeaver/ZFC/ConstructibleChoice.lean}.
\bibitem{danilkin2023}
Anna Danilkin and Lo\"ic Chevalier,
\emph{Three Squares Theorem}, Archive of Formal Proofs, 3 May 2023.
\url{https://isa-afp.org/entries/Three_Squares.html}.
\bibitem{sun2026}
Zhi-Wei Sun, \emph{Ten unknowns for Hilbert's tenth problem over the
integers}, arXiv:2609.01594v2, 2 September 2026.
\url{https://arxiv.org/abs/2609.01594}.
\end{thebibliography}
\end{document}
